\section{Related work}\label{sec:related}

Relabeling a fixed set of images along a class hierarchy is the experiment closest to ours.
Huh et al.\ merged the ImageNet classes along the WordNet tree and fine-tuned the resulting networks on detection and
scene classification, where the coarser labels cost little~\cite{huh}.
Hong et al.\ pretrained on levels of the class hierarchy of the full ImageNet dataset; their linear probes show a
larger gap between the finest and the coarsest labels than fine-tuning does, which they trace to the rarer features
that fine classes force a network to learn~\cite{hong}.
In weakly supervised pretraining the hashtag vocabulary that best matches the target classes transfers best, although
there the images change with the vocabulary~\cite{mahajan}.
The same contrast between frozen and adapted use appears when the training loss or the regularization changes instead
of the labels~\cite{kornblith,kornblithtransfer}, and fine-tuning can change the features it starts from enough to rank
representations differently~\cite{kumar}.
Training on finer labels can raise accuracy on the coarse classes themselves~\cite{chen}, though not on every
dataset~\cite{pirovano}, and representations finer than the training labels can be learned from coarse labels with
additional instance-level or nearest-neighbor objectives~\cite{grafit,xucoarse}.

What a representation keeps about merged classes has been studied more directly.
Late in training the last-layer features of each class contract to the class mean (neural collapse)~\cite{papyan},
a geometry argued both to help~\cite{galanti} and to hurt~\cite{hui} transfer, and one that would erase the subclasses
inside a merged class.
In practice subclasses often stay separable: networks trained on coarse labels leave their unlabeled subclasses
separable enough in feature space to be recovered by clustering~\cite{sohoni}, and a network made to divide each class
into subclasses of its own finds groupings close to the natural subclasses where these exist~\cite{subclassdistill}.
Yang, Steinhardt and Hu measured the effect of the merge itself: subclasses merged at random stay well separated in the
last-layer features, while subclasses merged into a natural superclass are distinguished less well~\cite{yang}.
Our random groupings follow theirs; we add transfer to other tasks, fine-tuning, and merges that each remove a named
physical distinction.

In jet physics, supervised classification of simulated jets is the pretraining objective of the Particle Transformer,
which showed that pretraining on the ten classes of JetClass improves top tagging and quark--gluon discrimination after
fine-tuning~\cite{part}, and of Sophon, which moved to the much finer classes of JetClass-II and to the direct use of
output scores and frozen features~\cite{sophon}.
CMS trains its taggers on many classes of decays~\cite{cmspart,glopartdp}, and supervised classification has been
compared with masked-particle and generative pretraining on the same jets, with the classes fixed~\cite{objectives}.
Self-supervised objectives~\cite{jetclr,rs3l,mpmtok,rieck,jbot,rino} such as masked particle modeling are evaluated
with a frozen and with a fine-tuned backbone as a function of the number of labeled jets~\cite{mpm}, the two modes we
also compare, and contrastive pretraining on growing subsets of JetClass has been fine-tuned on top tagging to measure
how the gain scales with the pretraining set~\cite{zhaossl}.
SuperJEPA combines a predictive self-supervised objective with classification and reports frozen-encoder and
fine-tuned results on JetClass and top tagging~\cite{superjepa}, and the features of a network trained for particle
reconstruction can be read out by a single linear layer for jet flavor identification~\cite{mlpf_fm}.
In these studies the classes of the supervised objective are fixed, or change together with the jets used for
pretraining.
